\documentclass[11pt]{article}
\usepackage[utf8]{inputenc}
\usepackage{amsmath, amssymb}
\usepackage{geometry}
\geometry{a4paper, margin=1in}
\usepackage{fancyhdr}
\usepackage{enumitem}
\usepackage{hyperref}

\pagestyle{fancy}
\fancyhead{}
\fancyhead[R]{DMS KU - AVP}
\fancyhead[L]{PDE Problem Sheet}
\fancyfoot[C]{\thepage}
\renewcommand{\headrulewidth}{0.4pt}

\newcounter{prob}
\newcommand{\prob}[1]{\refstepcounter{prob}\par\medskip\noindent\textbf{Problem \theprob.} \textit{(#1)}\quad}
\setlist{itemsep=2pt, topsep=3pt}

\begin{document}
\begin{center}
\textbf{\Large Partial Differential Equations} \\[0.5em]
\textbf{Course Code: MSMAT03DSC12} \\[0.5em]
\textbf{Problem Sheet: Heat Equation and Wave Equation} \\[0.5em]
\textbf{Anoop V. P.}
\end{center}

\noindent\textbf{Instructions.} These problems check your understanding of Section 3 (Heat Equation) and Section 4 (Wave Equation, including the Transport Equation) of the lecture notes. Each problem is a computation or a question that needs a short reason.

\smallskip
\noindent\textbf{Notation.} $\Phi(x,t) = (4\pi t)^{-n/2}e^{-|x|^2/4t}$ $(t > 0)$ is the fundamental solution of the heat equation. For the heat equation with initial value $g$, ``the solution'' means $u(x,t) = \int_{\mathbb{R}^n}\Phi(x-y,t)\,g(y)\,dy$. For the wave equation $u_{tt} - c^2u_{xx} = 0$ with $u(x,0) = g(x)$, $u_t(x,0) = h(x)$, ``the solution'' means the one given by d'Alembert's formula. You may use $\int_{-\infty}^\infty e^{-s^2}\,ds = \sqrt\pi$.

%=====================================================================
\section*{Part A: The Heat Equation}

\prob{Recognising solutions}
Decide which of the following functions solve the heat equation $u_t - \Delta u = 0$ in $\mathbb{R}^n \times (0,\infty)$ for the indicated $n$. For those that do not, find $f$ such that they solve $u_t - \Delta u = f$.
\begin{enumerate}[label=(\alph*)]
    \item $u(x,t) = x^2 + 2t$ \quad ($n = 1$).
    \item $u(x,t) = e^{-t}\sin x$ \quad ($n = 1$).
    \item $u(x,t) = e^{-k^2t}\cos(kx)$, $k \in \mathbb{R}$ \quad ($n = 1$).
    \item $u(x,t) = x^2 - 2t$ \quad ($n = 1$).
    \item $u(x,y,t) = x^2 + y^2 + 4t$ \quad ($n = 2$).
    \item $u(x,y,t) = e^{-2t}\sin x\sin y$ \quad ($n = 2$).
\end{enumerate}

\prob{Homogeneous and nonhomogeneous}
\begin{enumerate}[label=(\alph*)]
    \item Show that if $u_1, u_2$ solve the heat equation, then so does $au_1 + bu_2$ for constants $a, b$.
    \item Show that if $u_1$ solves $u_t - \Delta u = f$ and $u_2$ solves $u_t - \Delta u = 0$, then $u_1 + u_2$ solves $u_t - \Delta u = f$.
    \item Is the sum of two solutions of $u_t - \Delta u = f$ again a solution of $u_t - \Delta u = f$? Explain.
    \item Find one solution of $u_t - u_{xx} = 1$ that depends only on $t$, and one that depends only on $x$.
\end{enumerate}

\prob{The constant $a$}
Suppose $\mathbf F = -aDu$ with a constant $a > 0$ and $u_t = -\operatorname{div}\mathbf F$.
\begin{enumerate}[label=(\alph*)]
    \item Show that $u_t = a\Delta u$.
    \item Show that if $u$ solves $u_t = a\Delta u$, then $v(x,t) := u(x, t/a)$ solves $v_t = \Delta v$.
    \item Using (b) and $\Phi$, write down a solution of $u_t = a\Delta u$ in $\mathbb{R}^n \times (0,\infty)$.
\end{enumerate}

\prob{The scaling $u(\lambda x, \lambda^2t)$}
\begin{enumerate}[label=(\alph*)]
    \item Let $u(x,t) = x^2 + 2t$ $(n = 1)$ and $\lambda = 3$. Write down $w(x,t) = u(\lambda x, \lambda^2 t)$ and check directly that $w$ solves the heat equation.
    \item For the same $u$, show that $u(\lambda x, \lambda t)$ is \emph{not} a solution when $\lambda = 3$.
    \item Show that the ratio $\frac{|x|^2}{t}$ is unchanged under $(x,t) \mapsto (\lambda x, \lambda^2 t)$, but changes under $(x,t) \mapsto (\lambda x,\lambda t)$.
    \item Explain in two or three sentences why the scaling in (a) works but the one in (b) does not, in terms of the number of $t$- and $x$-derivatives in the equation.
\end{enumerate}

\prob{The dilation scaling}
Let $u$ solve the heat equation in $\mathbb{R}^n \times (0,\infty)$ and let $w(x,t) = \lambda^\alpha u(\lambda^\beta x, \lambda t)$ for a fixed $\lambda > 0$.
\begin{enumerate}[label=(\alph*)]
    \item Compute $w_t$ and $\Delta w$ in terms of derivatives of $u$.
    \item For which $\beta$ is $w$ a solution for every $\lambda > 0$? Does the answer depend on $\alpha$?
    \item Verify that $\Phi(x,t) = \lambda^{n/2}\,\Phi(\lambda^{1/2}x, \lambda t)$ for all $\lambda > 0$, $x \in \mathbb{R}^n$, $t > 0$.
    \item Setting $\lambda = t^{-1}$ in (c), express $\Phi(x,t)$ in the form $t^{-\alpha}v(x/t^\beta)$. What are $\alpha$, $\beta$ and $v$?
\end{enumerate}

\prob{Inserting the ansatz ($n = 1$)}
Let $n = 1$ and $u(x,t) = t^{-\alpha}v(y)$ with $y = t^{-\beta}x$.
\begin{enumerate}[label=(\alph*)]
    \item Compute $u_t$ and $u_{xx}$ in terms of $v(y)$, $v'(y)$, $v''(y)$ and powers of $t$.
    \item Write down the equation for $v$ that results from $u_t - u_{xx} = 0$, and show that $t$ drops out if and only if $\beta = \frac12$.
    \item With $\beta = \frac12$ and $\alpha = \frac12$, show that the equation for $v$ can be written as $\big(v' + \frac12 yv\big)' = 0$.
    \item Solve the equation in (c) under the assumption that $v$ and $v'$ tend to $0$ as $y \to \pm\infty$.
\end{enumerate}

\prob{The radial ODE}
\begin{enumerate}[label=(\alph*)]
    \item Let $v(y) = w(|y|)$ in $\mathbb{R}^n$ and $r = |y| \neq 0$. Show that $y\cdot Dv(y) = rw'(r)$.
    \item Write down the ODE for $w$ obtained from $\alpha v + \frac12 y\cdot Dv + \Delta v = 0$ when $n = 3$.
    \item Show that for $\alpha = \frac n2$, the ODE in (b) can be written as $(r^{n-1}w')' + \frac12(r^nw)' = 0$.
    \item Check directly that $w(r) = e^{-r^2/4}$ satisfies $w' = -\frac12 rw$ and hence the ODE in (c).
\end{enumerate}

\prob{$\Phi$ solves the heat equation}
\begin{enumerate}[label=(\alph*)]
    \item For $n = 1$, compute $\Phi_t$ and $\Phi_{xx}$ explicitly and check that $\Phi_t = \Phi_{xx}$ for $t > 0$.
    \item For $n = 2$, write $\Phi(x,t) = \Phi_1(x_1,t)\,\Phi_1(x_2,t)$ where $\Phi_1$ is the one-dimensional fundamental solution. Use this product form and (a) to show that $\Phi$ solves the heat equation in $\mathbb{R}^2 \times (0,\infty)$.
    \item Show that $(x,t) \mapsto \Phi(x - y, t - s)$ solves the heat equation for $t > s$, for each fixed $y \in \mathbb{R}^n$ and $s \in \mathbb{R}$.
\end{enumerate}

\prob{The shape of $\Phi$ ($n = 1$)}
Let $n = 1$.
\begin{enumerate}[label=(\alph*)]
    \item For fixed $t > 0$, where does $x \mapsto \Phi(x,t)$ attain its maximum, and what is the maximum value?
    \item Sketch $\Phi(\cdot,t)$ for $t = 1$, $t = \frac14$ and $t = \frac1{16}$ on the same axes. What happens to the height and to the width as $t$ decreases?
    \item Show that $\Phi(0,t) \to \infty$ and $\Phi(x,t) \to 0$ for each $x \neq 0$, as $t \to 0^+$.
    \item Fix $x \neq 0$. Find the time $t > 0$ at which $t \mapsto \Phi(x,t)$ is largest.
\end{enumerate}

\prob{Integrals of $\Phi$}
\begin{enumerate}[label=(\alph*)]
    \item For $n = 1$, show $\int_{-\infty}^\infty \Phi(x,t)\,dx = 1$ using the substitution $x = 2\sqrt t\,z$.
    \item For $n = 2$, compute $\int_{\mathbb{R}^2}\Phi(x,t)\,dx$ using polar coordinates.
    \item For $n = 1$, compute $\int_{-\infty}^\infty x^2\,\Phi(x,t)\,dx$. (You may use $\int_{-\infty}^\infty s^2e^{-s^2}ds = \frac{\sqrt\pi}{2}$.) How does it depend on $t$?
    \item For $n = 1$, show that $\int_{|x| \geq \delta}\Phi(x,t)\,dx = \frac{2}{\sqrt\pi}\int_{\delta/(2\sqrt t)}^\infty e^{-s^2}ds$, and deduce that it tends to $0$ as $t \to 0^+$ for each $\delta > 0$.
\end{enumerate}

\prob{The Gaussian integral}
\begin{enumerate}[label=(\alph*)]
    \item In $\mathbb{R}^2$, compute $\int_{B(0,R)} e^{-|x|^2}dx$ by polar coordinates, and let $R \to \infty$.
    \item Deduce $\int_{-\infty}^\infty e^{-s^2}ds = \sqrt\pi$, explaining where Fubini--Tonelli is used.
    \item Compute $\int_{\mathbb{R}^n} e^{-|x|^2}dx$ and $\int_{\mathbb{R}^n}e^{-a|x|^2}dx$ for $a > 0$.
\end{enumerate}

\prob{Solving initial-value problems}
In each case $n = 1$. Compute $u(x,t) = \int_{\mathbb{R}}\Phi(x - y,t)g(y)\,dy$ and verify that it solves the heat equation and $u(x,t) \to g(x)$ as $t \to 0^+$.
\begin{enumerate}[label=(\alph*)]
    \item $g(x) \equiv 5$.
    \item $g(x) = \cos x$. (You may use $\int_{-\infty}^\infty e^{-as^2}\cos(bs)\,ds = \sqrt{\pi/a}\;e^{-b^2/4a}$ for $a > 0$.)
    \item $g(x) = e^{-x^2}$. (Complete the square in the exponent.)
\end{enumerate}

\prob{Properties of the solution formula}
Let $g, g_1, g_2$ be bounded and continuous on $\mathbb{R}^n$, and let $u, u_1, u_2$ be the corresponding solutions.
\begin{enumerate}[label=(\alph*)]
    \item Show that $|u(x,t)| \leq \sup_{\mathbb{R}^n}|g|$ for all $x$ and $t > 0$.
    \item Show that if $g_1 \leq g_2$ everywhere, then $u_1 \leq u_2$ everywhere.
    \item Show that if $a \in \mathbb{R}^n$ and $\tilde g(x) = g(x - a)$, then the solution with initial value $\tilde g$ is $u(x - a, t)$.
    \item If $g$ is even ($n = 1$), show that $u(\cdot,t)$ is even for each $t > 0$.
\end{enumerate}

\prob{The hypotheses of the solution theorem}
\begin{enumerate}[label=(\alph*)]
    \item Let $g(x) = \min\{|x|, 1\}$ $(n = 1)$. Is $g$ differentiable? Is the solution $u$ differentiable in $\mathbb{R}\times(0,\infty)$? What is $\lim_{(x,t)\to(0,0),\,t>0}u(x,t)$?
    \item Let $g(x) = e^{x^2}$ $(n = 1)$. Show that $\int_{\mathbb{R}}\Phi(x - y,t)g(y)\,dy = \infty$ for $t \geq \frac14$. Which hypothesis of the theorem on the solution of the initial-value problem fails?
    \item In the proof that $u \in C^\infty(\mathbb{R}^n\times(0,\infty))$, explain why it is not enough to know that the derivatives of $\Phi$ are bounded on $\mathbb{R}^n \times [\delta,\infty)$ to differentiate under the integral sign over $\mathbb{R}^n$.
\end{enumerate}

\prob{Understanding the proof of the initial condition}
Let $g$ be bounded and continuous on $\mathbb{R}^n$ and $x^0 \in \mathbb{R}^n$.
\begin{enumerate}[label=(\alph*)]
    \item Show that $u(x,t) - g(x^0) = \int_{\mathbb{R}^n}\Phi(x - y,t)\big[g(y) - g(x^0)\big]\,dy$, and say which result of the notes is used.
    \item Show: if $|x - x^0| \leq \frac\delta2$ and $|y - x^0| \geq \delta$, then $|x - y| \geq \frac12|y - x^0|$. Draw a picture for $n = 1$.
    \item Why is it important in the proof that the estimate of the integral over $\{|y - x^0| \geq \delta\}$ does not depend on $x$?
    \item The proof gives $|u(x,t) - g(x^0)| < 2\varepsilon$ rather than $< \varepsilon$. Why is this enough?
\end{enumerate}

\prob{$\Phi$ and the Dirac delta}
Let $\varphi$ be bounded and continuous on $\mathbb{R}^n$.
\begin{enumerate}[label=(\alph*)]
    \item Show that $\int_{\mathbb{R}^n}\Phi(x,t)\varphi(x)\,dx = u(0,t)$, where $u$ is the solution with initial value $g = \varphi$. (Use that $\Phi(\cdot,t)$ is even.)
    \item Deduce that $\int_{\mathbb{R}^n}\Phi(x,t)\varphi(x)\,dx \to \varphi(0)$ as $t \to 0^+$.
    \item Explain how (b) and the fact that $\int_{\mathbb{R}^n}\Phi(x,t)\,dx = 1$ justify the notation ``$\Phi = \delta_0$ on $t = 0$''.
\end{enumerate}

\prob{Positivity}
Let $n = 1$ and $g(x) = \max\{0, 1 - |x|\}$.
\begin{enumerate}[label=(\alph*)]
    \item Sketch $g$. Is $g$ bounded and continuous? Where is $g = 0$?
    \item Show that the solution satisfies $u(100, 10^{-3}) > 0$.
    \item Show that $u(x,t) > 0$ for every $x \in \mathbb{R}$ and $t > 0$.
    \item Show that $u(x,t) \leq \frac{1}{\sqrt{4\pi t}}\int_{-1}^1 g(y)\,dy$, and deduce that $u(x,t) \to 0$ as $t \to \infty$, uniformly in $x$.
\end{enumerate}

%=====================================================================
\section*{Part B: The Transport Equation}

\prob{Solving homogeneous problems}
Solve the following initial-value problems and verify your answers.
\begin{enumerate}[label=(\alph*)]
    \item $u_t + 2u_x = 0$ in $\mathbb{R}\times(0,\infty)$, $u(x,0) = \sin x$.
    \item $u_t - 3u_x = 0$ in $\mathbb{R}\times(0,\infty)$, $u(x,0) = \dfrac{1}{1 + x^2}$.
    \item $u_t + u_{x_1} - u_{x_2} = 0$ in $\mathbb{R}^2\times(0,\infty)$, $u(x_1,x_2,0) = x_1x_2$.
    \item $u_t + b\cdot Du = 0$ in $\mathbb{R}^n\times(0,\infty)$, $u(x,0) = |x|^2$, for a given $b \in \mathbb{R}^n$.
\end{enumerate}

\prob{Geometry of the lines}
Consider $u_t + 2u_x = 0$ ($n = 1$, $b = 2$).
\begin{enumerate}[label=(\alph*)]
    \item In the $(x,t)$-plane, sketch the lines $L_y$ for $y = -2, 0, 2$. What is their slope?
    \item Find the line through the point $(5,1)$ and the point where it meets $t = 0$.
    \item Show that $(5,1)$ lies on no other line $L_y$.
    \item If $u(x,0) = g(x)$ and $u(5,1) = 7$, what can you say about $g$? What is $u(3 + 2s, s)$ for $s \geq 0$?
\end{enumerate}

\prob{Lines in higher dimensions}
Let $n = 2$ and $b = (1,2)$.
\begin{enumerate}[label=(\alph*)]
    \item Describe the half-line $L_y$ starting at $y = (0,0)$ in $\mathbb{R}^2 \times [0,\infty)$.
    \item Which line $L_y$ passes through the point $(x,t) = ((3,4),1)$?
    \item Show that the map $(y,\sigma) \mapsto (y + \sigma b, \sigma)$ is a bijection of $\mathbb{R}^2\times[0,\infty)$ onto itself, and find its inverse.
    \item Explain, using (c), why the lines $L_y$ are pairwise disjoint and fill $\mathbb{R}^2\times[0,\infty)$.
\end{enumerate}

\prob{Transport of a profile}
Let $u$ solve $u_t + bu_x = 0$ ($n = 1$), $u(x,0) = e^{-x^2}$, with $b > 0$.
\begin{enumerate}[label=(\alph*)]
    \item Write down $u$. Where does $u(\cdot,t)$ attain its maximum, and what is the maximum value?
    \item Sketch $u(\cdot,t)$ for $t = 0, 1, 2$. Describe what happens to the graph.
    \item Compare with the heat equation with the same initial value (Problem 12(c)): what happens there to the maximum value as $t$ increases?
\end{enumerate}

\prob{Solving nonhomogeneous problems}
Solve the following problems using the solution formula for the nonhomogeneous transport equation, and verify your answers.
\begin{enumerate}[label=(\alph*)]
    \item $u_t + u_x = t$, $u(x,0) = 0$.
    \item $u_t + 2u_x = x$, $u(x,0) = 0$.
    \item $u_t - u_x = e^t$, $u(x,0) = \cos x$.
    \item $u_t + u_x = \sin t$, $u(x,0) = x^2$.
\end{enumerate}

\prob{A transport equation with a zero-order term}
Let $c \in \mathbb{R}$ and $b \in \mathbb{R}^n$, and consider $u_t + b\cdot Du + cu = 0$ in $\mathbb{R}^n\times(0,\infty)$, $u = g$ on $t = 0$.
\begin{enumerate}[label=(\alph*)]
    \item Show that $v(x,t) := e^{ct}u(x,t)$ solves $v_t + b\cdot Dv = 0$.
    \item Find $v(x,0)$ and solve the problem for $u$.
    \item Alternatively, compute $\dot z(s)$ for $z(s) = u(x + sb, t + s)$, solve the resulting ODE, and obtain the same formula.
\end{enumerate}

\prob{The role of the hypotheses}
\begin{enumerate}[label=(\alph*)]
    \item Let $g(x) = \sin x$ and define $u(x,t) = g(x - t)$ for $t > 0$ and $u(x,0) = 0$. Show that $u$ satisfies $u_t + u_x = 0$ for $t > 0$ and $u(x,0) = 0$, but $u \not\equiv 0$. Which requirement in the definition of a solution fails?
    \item Let $g(x) = |x|$ and $b = 1$. At which points $(x,t)$ is $u(x,t) = g(x - t)$ not differentiable? Sketch this set.
    \item Show that if $u \in C^1(\mathbb{R}\times(0,\infty))$ solves $u_t + u_x = 0$ and $u(x,1) = |x|$ for all $x$, we get a contradiction.
\end{enumerate}

%=====================================================================
\section*{Part C: The Wave Equation}

\prob{Recognising solutions}
\begin{enumerate}[label=(\alph*)]
    \item Decide which of the following solve $u_{tt} - u_{xx} = 0$: \ $\sin x\cos t$, \ $(x + t)^3$, \ $e^{x - t}$, \ $xt$, \ $x^2 + t^2$, \ $x^2 - t^2$, \ $\sin x\sin 2t$.
    \item For each solution in (a), write it in the form $F(x + t) + G(x - t)$.
    \item Show that both $u \equiv 0$ and $u = t$ solve $u_{tt} - u_{xx} = 0$ with $u(x,0) = 0$. What does this say about the number of initial conditions needed?
    \item For which constants $k$ does $u = \sin(x)\cos(kt)$ solve $u_{tt} - c^2u_{xx} = 0$?
\end{enumerate}

\prob{A radial wave in three dimensions}
Let $n = 3$, $r = |x| > 0$, $F \in C^2(\mathbb{R})$, and $u(x,t) = \dfrac{F(r - t)}{r}$.
\begin{enumerate}[label=(\alph*)]
    \item Using the radial formula $\Delta u = u_{rr} + \frac{n-1}{r}u_r$, show that $\Delta u = \dfrac{F''(r - t)}{r}$.
    \item Show that $u_{tt} - \Delta u = 0$ for $x \neq 0$.
    \item Is $\dfrac{F(r + t)}{r}$ also a solution?
\end{enumerate}

\prob{The factorization}
\begin{enumerate}[label=(\alph*)]
    \item For $u \in C^2$, verify $\big(\partial_t + c\,\partial_x\big)\big(\partial_t - c\,\partial_x\big)u = u_{tt} - c^2u_{xx}$. Where is $u \in C^2$ used?
    \item Show that the two factors commute: $\big(\partial_t - c\,\partial_x\big)\big(\partial_t + c\,\partial_x\big)u = u_{tt} - c^2u_{xx}$ as well.
    \item Show that every $C^2$ solution of $u_t - cu_x = 0$ and every $C^2$ solution of $u_t + cu_x = 0$ solves the wave equation $u_{tt} - c^2u_{xx} = 0$.
    \item Which transport equation is solved by $F(x + ct)$, and which by $G(x - ct)$?
\end{enumerate}

\prob{The derivation, step by step}
Let $u(x,t) = x^2 + t^2$, the solution of $u_{tt} - u_{xx} = 0$ with $g(x) = x^2$, $h(x) = 0$.
\begin{enumerate}[label=(\alph*)]
    \item Compute $v = u_t - u_x$ and $a(x) = v(x,0)$. Check that $a = h - g'$.
    \item Check that $v(x,t) = a(x - t)$, and that $v_t + v_x = 0$. Which transport equation (which $b$) is this?
    \item Check that $u$ solves $u_t - u_x = a(x - t)$. Which transport equation (which $b$, which $f$) is this?
    \item Recover $u$ from the formula $u(x,t) = g(x + t) + \frac12\int_{x-t}^{x+t}a(y)\,dy$.
\end{enumerate}

\prob{d'Alembert's formula}
Solve $u_{tt} - u_{xx} = 0$, $u(x,0) = g(x)$, $u_t(x,0) = h(x)$ for:
\begin{enumerate}[label=(\alph*)]
    \item $g(x) = \sin x$, $h(x) = 0$;
    \item $g(x) = 0$, $h(x) = \cos x$;
    \item $g(x) = x^2$, $h(x) = x$;
    \item $g(x) = e^{-x^2}$, $h(x) = 0$.
\end{enumerate}
In (a)--(c) simplify the answer and verify it. In (d) sketch $u(\cdot,t)$ for $t = 0$, $t = 1$ and $t = 4$.

\prob{The wave equation with speed $c$}
\begin{enumerate}[label=(\alph*)]
    \item Solve $u_{tt} = 4u_{xx}$, $u(x,0) = \sin x$, $u_t(x,0) = 2\cos x$, and simplify. Is the answer of the form $F(x + 2t)$, $G(x - 2t)$, or a sum of both?
    \item Solve $u_{tt} = 4u_{xx}$, $u(x,0) = 0$, $u_t(x,0) = \dfrac{1}{1 + x^2}$.
    \item Solve $u_{tt} = 9u_{xx}$, $u(x,0) = x$, $u_t(x,0) = x^2$.
    \item Let $u$ solve $u_{tt} = c^2u_{xx}$. Show that $w(x,\tau) = u(x,\tau/c)$ solves $w_{\tau\tau} = w_{xx}$, and find the initial values of $w$ in terms of $g$, $h$, $c$.
\end{enumerate}

\prob{General solution and initial values}
Suppose $u(x,t) = F(x + t) + G(x - t)$ with $F, G \in C^2(\mathbb{R})$, and $u(x,0) = g(x)$, $u_t(x,0) = h(x)$.
\begin{enumerate}[label=(\alph*)]
    \item Show that $F + G = g$ and $F' - G' = h$.
    \item Deduce $F - G = \int_0^x h(y)\,dy + K$ for a constant $K$, and solve for $F$ and $G$.
    \item Substitute into $u$ and recover d'Alembert's formula. What happens to $K$?
    \item Use d'Alembert's formula to show that if $g \equiv 0$ and $h \equiv 0$, then the only $C^2$ solution is $u \equiv 0$. Deduce that two $C^2$ solutions with the same $g, h$ coincide.
\end{enumerate}

\prob{No smoothing}
Let $g(x) = x|x|$ and $h(x) = 0$.
\begin{enumerate}[label=(\alph*)]
    \item Show that $g \in C^1(\mathbb{R})$ but $g \notin C^2(\mathbb{R})$.
    \item Write down $u$ from d'Alembert's formula. Show that $u_{tt} = u_{xx}$ except on the lines $x = t$ and $x = -t$.
    \item On which lines is $u$ not $C^2$? Sketch them. Is $u(\cdot,t)$ smoother than $g$ for $t > 0$?
    \item Compare with the heat equation: is the solution with a bounded, continuous, non-differentiable initial value (such as $\min\{|x|,1\}$) smooth for $t > 0$?
\end{enumerate}

\prob{Odd and even data}
Let $g \in C^2(\mathbb{R})$, $h \in C^1(\mathbb{R})$ and let $u$ solve $u_{tt} - u_{xx} = 0$, $u = g$, $u_t = h$ at $t = 0$.
\begin{enumerate}[label=(\alph*)]
    \item If $g$ and $h$ are odd, show that $u(0,t) = 0$ for all $t \geq 0$.
    \item If $g$ and $h$ are even, show that $u(\cdot,t)$ is even and $u_x(0,t) = 0$ for all $t \ge 0$.
    \item If $g$ and $h$ are $2\pi$-periodic, show that $u(\cdot,t)$ is $2\pi$-periodic for each $t$.
\end{enumerate}

\prob{Domains of dependence and influence}
Consider $u_{tt} - 4u_{xx} = 0$ ($c = 2$).
\begin{enumerate}[label=(\alph*)]
    \item Find the domain of dependence of the point $(1,3)$.
    \item Find the range of influence of the point $(0,0)$, and the segment it occupies at time $t = 1$.
    \item Find the range of influence of the interval $[-1,1]$, and the segment it occupies at time $t = 2$.
    \item Find the domain of determinacy of $[0,4]$: its vertices, its height, and the inequalities describing it.
    \item Decide whether the point $(3,1)$ lies in each of the sets in (b), (c), (d).
\end{enumerate}

\prob{Data with compact support}
Let $c = 1$, $g \equiv 0$ and $h(x) = (1 - x^2)^2$ for $|x| \leq 1$, $h(x) = 0$ for $|x| > 1$.
\begin{enumerate}[label=(\alph*)]
    \item Show that $h \in C^1(\mathbb{R})$.
    \item For $t = 2$, for which $x$ is $u(x,2) = 0$? Relate your answer to the range of influence of $[-1,1]$.
    \item Show that $u(x,t) = \frac{8}{15}$ whenever $|x| \leq t - 1$.
    \item For the heat equation with the same initial value $h$, is the solution zero anywhere for $t > 0$? Explain the difference.
\end{enumerate}

\prob{Domain of determinacy}
Let $c = 1$ and $[a,b] = [0,2]$.
\begin{enumerate}[label=(\alph*)]
    \item Sketch the domain of determinacy $T$ of $[0,2]$ and the range of influence of $[0,2]$ on the same picture.
    \item Show that if two sets of data $(g_1,h_1)$, $(g_2,h_2)$ agree on $[0,2]$, then $u_1(1,\frac12) = u_2(1,\frac12)$.
    \item The point $(2,1)$ is not in $T$. Give explicit data $(g_1,h_1)$ and $(g_2,h_2)$, both $C^1$ in $h$ and $C^2$ in $g$, which agree on $[0,2]$ but give $u_1(2,1) \neq u_2(2,1)$.
    \item For a point $(x,t)$ in the range of influence of $[0,2]$ but not in $T$, is $u(x,t)$ determined by the data on $[0,2]$ alone? Explain.
\end{enumerate}

\prob{Heat versus wave}
Answer each question with one or two sentences, giving the reason from the notes.
\begin{enumerate}[label=(\alph*)]
    \item How many initial conditions are prescribed for the heat equation, and how many for the wave equation? Why?
    \item If the initial data are continuous but not differentiable, is the solution smooth for $t > 0$? (Answer for both equations.)
    \item If the initial data vanish outside $[-1,1]$, is the solution at time $t = 1$ zero outside some bounded interval? (Answer for both equations.)
    \item Does $u(x,t)$ depend on the initial data on all of $\mathbb{R}$, or only on part of it? (Answer for both equations, $n = 1$.)
\end{enumerate}

\end{document}
